%% RevTeX4-2 document
\documentclass[aps,prb,twocolumn,superscriptaddress,floatfix]{revtex4-2}

%% ── Packages ────────────────────────────────────────────────────────────────
\usepackage{amsmath,amssymb,amsfonts}
\usepackage{mathtools}
\usepackage{bm}
\usepackage{graphicx}
\usepackage{xcolor}
\usepackage{booktabs}
\usepackage{hyperref}
%% ── Macros ──────────────────────────────────────────────────────────────────
\newcommand{\Nx}{N_x}
\newcommand{\Ny}{N_y}
\newcommand{\hmod}{K_A}
\newcommand{\GA}{G_A}
\newcommand{\GAtop}{G^{\mathrm{top}}}
\newcommand{\half}{\tfrac{1}{2}}
\newcommand{\C}{\mathcal{C}}
\newcommand{\Htop}{H_{\mathrm{CI}}}
\newcommand{\Hflat}{H_{\mathrm{flat}}}
\newcommand{\Nxy}{N_{xy}}
\DeclareMathOperator{\arctanh}{arctanh}

\begin{document}

\title{Modular-Hamiltonian Charge Spreading in the\\
       Class-A Fermionic Adaptive Circuit:\\
       Exact Ground State vs.\ Overcomplete-Wannier Surrogates}

\date{\today}
\maketitle

%% ─────────────────────────────────────────────────────────────────────────────
\section{Overview}
\label{sec:overview}

This note documents a diagnostic computation that probes the half-system
modular Hamiltonian of the domain-wall (DW) Chern-insulator ground state and
a family of overcomplete-Wannier (OW) flattened-$\sigma_x$ surrogate states.
The central question is whether the OW surrogates—which are used as the
adaptive-measurement target in the circuit—faithfully capture the
entanglement-theoretic structure encoded in the exact many-body ground state,
as revealed by modular-time charge propagation.

We work on an $\Nx \times \Ny = 20 \times 40$ torus with a domain-wall mass
profile, build the modular Hamiltonian from the half-cylinder averaged
covariance matrix, inject unit charges at the domain-wall sites, and evolve
them under the resulting quadratic generator.  The key observables are the
charge-density heatmap $\Nxy(t)$ and the fraction of charge retained inside
the topological slab.

%% ─────────────────────────────────────────────────────────────────────────────
\section{Model and Geometry}
\label{sec:model}

\subsection{Class-A Chern insulator with a domain wall}

We study a two-band lattice Dirac model on a rectangular $\Nx \times \Ny$
torus (periodic in both directions).  The single-particle Hamiltonian in
real space is
\begin{equation}
  \Htop = \sum_{\mathbf{r},\mathbf{r}'} c^\dagger_{\mathbf{r}}
           h_{\mathbf{r}\mathbf{r}'} c_{\mathbf{r}'},
\end{equation}
with the $2\times 2$ Bloch kernel (written at wavevector $\mathbf{k}$
for the translationally symmetric case)
\begin{equation}
  h(\mathbf{k}) = \hat{n}(\mathbf{k}) \cdot \bm{\sigma},
  \label{eq:bloch}
\end{equation}
where $\bm{\sigma} = (\sigma_x, \sigma_y, \sigma_z)$ are Pauli matrices and
\begin{align}
  n_x(\mathbf{k}) &= \sin k_x, \\
  n_y(\mathbf{k}) &= \sin k_y, \\
  n_z(x,\mathbf{k}) &= \alpha(x) - \cos k_x - \cos k_y.
\end{align}
The spatially varying mass $\alpha(x)$ defines a domain-wall profile:
$\alpha(x) = \alpha_1 = 1$ for $x \in [6, 14]$ (topological slab,
$C = +1$) and $\alpha(x) = \alpha_2 = 30$ for $x \notin [6, 14]$
(trivial region, $C = 0$).  The DW locations at $x = 6$ and $x = 14$
host one chiral edge mode each. The real-space hoppings are
\begin{align}
  h^x_{\mathrm{hop}} &= -\half \sigma_z - \frac{i}{2}\sigma_x, \\
  h^y_{\mathrm{hop}} &= -\half \sigma_z - \frac{i}{2}\sigma_y,
\end{align}
with the diagonal onsite term $\alpha(x)\sigma_z$.

\subsection{Covariance matrix convention}

All computations use the Majorana-convention \emph{reduced covariance matrix}
\begin{equation}
  G_{ij} = \langle c_i c^\dagger_j \rangle - \langle c^\dagger_j c_i\rangle
          = 2\langle c^\dagger_j c_i\rangle - \delta_{ij},
\end{equation}
so that $G = -G^\dagger$ and $G^2 = -\mathbf{1}$ for a pure Gaussian state.
The relation to the density-matrix covariance is $C = (\mathbf{1} + G)/2$.
For the exact DW ground state $G$ is obtained by diagonalizing $\Htop$ and
filling all negative-energy modes.

%% ─────────────────────────────────────────────────────────────────────────────
\section{Overcomplete-Wannier Surrogate States}
\label{sec:ow}

\subsection{Flattened $\sigma_x$ Hamiltonian}

The adaptive circuit uses a family of OW projectors built from the
band projectors of \eqref{eq:bloch}.  Given trial spinors
$\tau_A = \frac{1}{\sqrt{2}}(1,1)^T$ and
$\tau_B = \frac{1}{\sqrt{2}}(1,-1)^T$ (eigenstates of $\sigma_x$),
the projected Wannier amplitudes centered at $\mathbf{R}$ are
\begin{equation}
  W^{\pm}_{A(B),\mathbf{R}}(\mathbf{r}) =
    \sum_{\mathbf{k}} e^{i\mathbf{k}\cdot(\mathbf{r}-\mathbf{R})}
    P_\pm(\mathbf{k})\,\tau_{A(B)},
\end{equation}
with $P_\pm(\mathbf{k}) = \half(\mathbf{1} \pm \hat{n}\cdot\bm{\sigma})$ the
positive/negative band projectors.  Each Wannier mode is optionally truncated
to a shell of $n_\mathrm{sh}$ lattice sites around its center (the
\texttt{nshell} parameter) and renormalized, keeping only the $n_\mathrm{sh}$
nearest-neighbor ring in the $(x,y)$ plane.

The \emph{flattened} many-body state is the ground state of
\begin{equation}
  \Hflat = \sum_{\mathbf{R}}
           \bigl(
             W^+_{A,\mathbf{R}} (W^+_{A,\mathbf{R}})^\dagger
           + W^+_{B,\mathbf{R}} (W^+_{B,\mathbf{R}})^\dagger
           - W^-_{A,\mathbf{R}} (W^-_{A,\mathbf{R}})^\dagger
           - W^-_{B,\mathbf{R}} (W^-_{B,\mathbf{R}})^\dagger
           \bigr),
  \label{eq:Hflat}
\end{equation}
and the covariance $G_{\mathrm{flat}}$ is built from its occupied eigenstates.

\subsection{Variant constructions}

We study four variants for each $n_\mathrm{sh} \in \{1, 2\}$, differing in
how the trivial region is treated:

\begin{description}
  \item[\textbf{Standard}] No extra spatial truncation beyond \texttt{nshell}.
  \item[\textbf{DW-truncation}]  After the $n_\mathrm{sh}$ cutoff each
    Wannier mode is additionally masked to have support only within the
    region—topological or trivial—that contains its Wannier center.  This
    prevents inter-region tunneling in $\Hflat$ at the cost of breaking
    translation invariance near the DW.
  \item[\textbf{Trivial-region local mode}]  In the trivial region the
    real-space CI Hamiltonian is replaced by decoupled onsite blocks $-\mathbf{1}_2$,
    freezing the trivial-region orbitals into local maximally-filled modes.
  \item[\textbf{DW-truncation + trivial local}]  Both modifications applied
    simultaneously.
\end{description}

%% ─────────────────────────────────────────────────────────────────────────────
\section{Modular Hamiltonian Construction}
\label{sec:modular}

\subsection{Half-cylinder partition and $y$-averaging}

Fix a bipartition of the cylinder into the subsystem
$A(y_0) = [0, \Nx) \times [y_0, y_0 + \Ny/2)$ (in coordinates modulo $\Ny$)
and its complement $\bar{A}(y_0)$.  For a Gaussian state with covariance $G$
the restricted covariance is
\begin{equation}
  G_A(y_0) = G\big|_{A(y_0) \times A(y_0)}.
\end{equation}
Because the system is translation-invariant along $y$ (at fixed $x$),
physical observables are independent of $y_0$.  We therefore
\emph{average} over all $\Ny$ translations before extracting the modular
Hamiltonian:
\begin{equation}
  \overline{G}_A = \frac{1}{\Ny} \sum_{y_0=0}^{\Ny-1}
      U_{y_0}\, G_A(y_0)\, U_{y_0}^\dagger,
  \label{eq:Gavg}
\end{equation}
where $U_{y_0}$ is the unitary that relabels $A(y_0)$ into a common relative-$y$
frame $[0,\Ny/2) \times [0,\Nx)$.  The averaged matrix $\overline{G}_A$ is
Hermitian and we symmetrize numerically.  The dimension of $\overline{G}_A$ is
$\Nx \cdot (\Ny/2) \cdot 2 = 800$.

\subsection{Modular Hamiltonian}

The modular (entanglement) Hamiltonian $\hmod$ is defined by
\begin{equation}
  \rho_A \propto e^{-\hmod},
\end{equation}
and for a Gaussian state it is quadratic.  Writing the spectral decomposition
$\overline{G}_A = V \operatorname{diag}(g_i) V^\dagger$ with
$g_i \in (-1, 1)$, the modular Hamiltonian kernel is
\begin{equation}
  \hmod = -2\, V \operatorname{diag}\!\bigl(\arctanh g_i\bigr) V^\dagger.
  \label{eq:hmod}
\end{equation}
Equation~\eqref{eq:hmod} follows from the relation between the reduced density
matrix of a Gaussian state and its single-particle covariance \cite{Peschel2003}:
\begin{equation}
  \rho_A = \frac{1}{Z}
    \exp\!\left(
      -\sum_{ij} [h_{\mathrm{mod}}]_{ij}\, c^\dagger_i c_j
    \right),
    \quad
  h_{\mathrm{mod}} = -\ln \frac{C_A}{\mathbf{1}-C_A},
\end{equation}
where $C_A = (\mathbf{1} + \overline{G}_A)/2$.  The eigenvalues $\varepsilon_i =
-2\arctanh(g_i)$ are the single-particle modular energies; they span
$[-|\varepsilon|_{\max}, |\varepsilon|_{\max}]$ symmetrically about zero for a
particle-hole symmetric state.

%% ─────────────────────────────────────────────────────────────────────────────
\section{Charge-Injection Protocol}
\label{sec:injection}

\subsection{Initial state}

We inject charge by constructing an initial single-particle wavefunction
$\phi_0$ that places exactly one unit of charge in \emph{each orbital}
$(\mu = 0, 1)$ of every selected unit cell $(x, y_\mathrm{rel})$:
\begin{equation}
  \phi_0 = \sum_{(x,y,\mu) \in \mathcal{S}} \ket{x, y_\mathrm{rel}, \mu},
\end{equation}
where $\mathcal{S}$ is the injection set.  We use two injection
configurations, both evaluated at the reduced $y$ values $\{0, \Ny/2 - 1\}$
(the two edges of the half-system):
\begin{description}
  \item[\textbf{Config~1}] $x \in \{6, 14\}$ (DW sites only),
    $|\mathcal{S}| = 4$, total charge $Q_0 = 8$.
  \item[\textbf{Config~2}] $x \in \{5, 6, 7, 13, 14, 15\}$ (DW $\pm 1$
    neighborhood), $|\mathcal{S}| = 12$, total charge $Q_0 = 24$.
\end{description}

\subsection{Exact spectral evolution}

The wavefunction is evolved under $\hmod$ in the eigenbasis of $\overline{G}_A$.
Expanding $\phi_0$ in the modular eigenstates $\{v_i\}$,
\begin{equation}
  c_i = v_i^\dagger \phi_0,
\end{equation}
the time-evolved wavefunction is
\begin{equation}
  \phi(t) = \sum_i c_i\, e^{-i\varepsilon_i t}\, v_i,
  \label{eq:phit}
\end{equation}
and the charge density is
\begin{equation}
  \Nxy(t) = \sum_\mu |\phi_{x,y,\mu}(t)|^2.
\end{equation}
The time step is $\delta t = 0.05$, the total evolution time is
$T = 2\Nx = 40$, giving 800 steps.  Both the total charge $Q(t) = \sum_{xy}
\Nxy(t)$ and the norm $\|\phi(t)\|^2$ are conserved exactly by
\eqref{eq:phit}; numerically these drift by $< 5 \times 10^{-7}$ and
$< 3 \times 10^{-14}$, respectively—consistent with floating-point
round-off in the matrix-vector products.

%% ─────────────────────────────────────────────────────────────────────────────
\section{Results}
\label{sec:results}

\subsection{Universality of the modular spectrum}

\begin{table}[t]
\caption{Half-system modular-Hamiltonian spectrum summary.
  All 9 state variants share an identical spectrum: 800 modes,
  $\varepsilon_{\min} = -23.719$, $\varepsilon_{\max} = +23.719$,
  Hermiticity error $= 0$, averaged over all $y_0$ half-cuts.}
\label{tab:spectra}
\begin{ruledtabular}
\begin{tabular}{lcccc}
  State & $n_\mathrm{sh}$ & $\varepsilon_{\min}$ & $\varepsilon_{\max}$
        & $\|K - K^\dagger\|_\infty$ \\
  \hline
  Exact DW GS             & --  & $-23.719$ & $+23.719$ & $0$ \\
  OW standard             &  1  & $-23.719$ & $+23.719$ & $0$ \\
  OW standard             &  2  & $-23.719$ & $+23.719$ & $0$ \\
  OW triv-local           &  1  & $-23.719$ & $+23.719$ & $0$ \\
  OW triv-local           &  2  & $-23.719$ & $+23.719$ & $0$ \\
  OW DW-trunc             &  1  & $-23.719$ & $+23.719$ & $0$ \\
  OW DW-trunc             &  2  & $-23.719$ & $+23.719$ & $0$ \\
  OW DW-trunc + triv-local&  1  & $-23.719$ & $+23.719$ & $0$ \\
  OW DW-trunc + triv-local&  2  & $-23.719$ & $+23.719$ & $0$ \\
\end{tabular}
\end{ruledtabular}
\end{table}

The most striking result is that \emph{every} state variant—from the exact
DW ground state to the most aggressively truncated OW surrogate—produces an
identical modular Hamiltonian spectrum (Table~\ref{tab:spectra}).  The 800
single-particle modular energies span $[-23.719, +23.719]$ symmetrically,
with machine-precision Hermiticity.  This universality arises because the
$y$-averaged half-system covariance $\overline{G}_A$ in \eqref{eq:Gavg} is
governed by the long-wavelength entanglement structure of the topological
phase—dominated by the $\Ny$ chiral edge modes crossing the cut—rather than
by the fine-grained real-space details of the wavefunction.  The fact that
all OW surrogates reproduce this spectrum exactly confirms that they encode
the same topological entanglement as the exact ground state.

\subsection{Charge-density dynamics: snapshots}

\begin{figure}[h]
  \includegraphics[width=\columnwidth]{fig_snapshots.png}
  \caption{Charge-density heatmaps $\Nxy(t)$ at $t = 0, T/4, T/2, T$
    for three representative states (Config~1, injection at DW sites
    $x \in \{6,14\}$, edge rows $y \in \{0, 19\}$).
    Red dashed lines mark the domain-wall boundaries.
    Top row: exact DW ground state.  Middle: OW $n_\mathrm{sh}=1$.
    Bottom: OW $n_\mathrm{sh}=1$ with DW-truncation.  The exact and
    standard OW states show charge leaking into the trivial region at
    long times, while DW-truncation strictly confines the spreading.}
  \label{fig:snapshots}
\end{figure}

Figure~\ref{fig:snapshots} shows the space-time structure of the charge.
At $t=0$ all charge is localized at the injection sites at the two DW
columns ($x=6,14$) and the two $y$-edges.  Under modular-time evolution
the charge spreads primarily \emph{along} the domain walls (the $y$
direction), consistent with the chiral nature of the edge mode.

At intermediate times ($T/4$, $T/2$) the charge fans out into the bulk of
the topological slab for all three variants, reflecting the finite bandwidth
of the modular eigenstates.  The differences between variants become
pronounced at $T$: the exact and standard OW states allow a non-negligible
fraction of charge to escape into the trivial region ($x < 6$ or $x > 14$),
while the DW-truncation variant keeps all charge strictly inside the slab.

\subsection{$x$-marginal dynamics}

\begin{figure}[h]
  \includegraphics[width=\columnwidth]{fig_xmarginal.png}
  \caption{$x$-marginal charge $\sum_y \Nxy(t)$ as a function of
    modular time, shown as a two-dimensional heatmap (color intensity
    proportional to charge at each $x$).  Red dashed lines at $x = 6$
    and $x = 14$ mark the domain-wall positions.  The DW-truncation
    variant (right panel) shows absolute confinement to $x \in [6, 14]$
    at all times.}
  \label{fig:xmarg}
\end{figure}

Figure~\ref{fig:xmarg} shows the $x$-marginal charge $\sum_y \Nxy(t)$
as a space-time heatmap.  The exact DW ground state develops a broad,
roughly uniform distribution within the topological slab by $t \sim T/2$,
with a small exponentially suppressed tail that grows slowly into the
trivial region.  The standard OW surrogates behave qualitatively similarly
but with modestly enhanced spread.  The DW-truncation variant shows a
strictly binary spatial structure: charge is exactly zero outside the
slab for all $t$.

\subsection{Topological charge fraction}

\begin{figure}[h]
  \includegraphics[width=\columnwidth]{fig_top_fraction.png}
  \caption{Fraction of total charge $Q_\mathrm{top}(t)/Q_\mathrm{tot}$
    retained in the topological slab $x \in [6, 14]$ as a function of
    modular time $t$ (Config~1).  DW-truncation variants (dotted lines)
    maintain the fraction at exactly unity for all $t$.  Standard OW
    surrogates (dashed) relax towards a lower asymptote, following the
    exact DW GS (solid black) closely at short times.}
  \label{fig:topfrac}
\end{figure}

Figure~\ref{fig:topfrac} and Table~\ref{tab:topfrac} quantify the charge
leakage.  At $t = T = 40$ the exact DW ground state retains $95.9\%$ of
charge inside the topological slab.  The OW $n_\mathrm{sh}=1$ and
$n_\mathrm{sh}=2$ standards retain $89.6\%$ and $85.9\%$, respectively.
Both DW-truncation variants retain $100\%$ at all times.

\begin{table}[h]
\caption{Topological charge fraction $Q_\mathrm{top}(T)/Q_\mathrm{tot}$
  at $t = T = 40$ for Config~1.}
\label{tab:topfrac}
\begin{ruledtabular}
\begin{tabular}{lccc}
  State & $n_\mathrm{sh}$ & $Q_\mathrm{top}(T)/Q_\mathrm{tot}$ \\
  \hline
  Exact DW GS         &  -- & $0.959$ \\
  OW standard         &   1 & $0.896$ \\
  OW standard         &   2 & $0.859$ \\
  OW DW-truncation    &   1 & $1.000$ \\
  OW DW-truncation    &   2 & $1.000$ \\
\end{tabular}
\end{ruledtabular}
\end{table}

The residual leakage of the exact ground state ($\approx 4\%$) is a
finite-size effect: the modular Hamiltonian of the exact DW GS is
not exactly block-diagonal in the topological/trivial partition because
the Wannier functions of the occupied bands have exponentially decaying
but nonzero tails into the trivial region.  The slightly larger leakage of
the standard OW surrogates indicates that the $y$-averaged modular
eigenstates have a somewhat wider real-space support—a consequence of the
Wannier truncation distorting the long-range structure of the covariance.

%% ─────────────────────────────────────────────────────────────────────────────
\section{Physical Interpretation and Implications}
\label{sec:discussion}

\subsection{Modular time as an entanglement probe}

The modular (entanglement) Hamiltonian $\hmod$ generates a one-parameter
unitary group $e^{-i\hmod t}$ that acts purely within the subsystem $A$.
Evolving a local excitation under $\hmod$ traces out a trajectory in the
Hilbert space of $A$ that is entirely determined by the reduced density
matrix $\rho_A$—and hence by the entanglement structure of the underlying
state.  The speed and geometry of charge spreading under $\hmod$ is therefore
a direct indicator of the pattern of bipartite correlations.

For a topological state the entanglement Hamiltonian has a distinctive
\emph{chiral edge} structure: its low-energy spectrum is dominated by modes
localized near the entanglement cut (the $y$-boundaries of $A$), moving
in one direction along the cut.  These modes are the entanglement analogue
of the physical chiral edge modes of the Chern insulator.  Charge injected
at the DW columns—where topological edge modes are concentrated—is
efficiently coupled to these entanglement-edge modes and spreads along $y$.
Charge that attempts to spread along $x$ into the trivial region must
overcome a large modular ``gap'' (the trivial-region modes have
$\varepsilon \to \pm \infty$ in the $\alpha_2 \gg 1$ limit), leading to
the observed quasi-confinement.

\subsection{Universality of the modular spectrum}

The exact equality of all nine modular spectra in Table~\ref{tab:spectra}
is not accidental.  The $y$-averaging procedure in \eqref{eq:Gavg} projects
$\overline{G}_A$ onto the translation-invariant sector along $y$.  Since the
CI is gapped in the $x$ bulk and the DW hosts topologically protected gapless
edge modes, the dominant entanglement entropy is carried by these edge modes.
In the large-$\alpha_2$ limit the OW Wannier functions in the trivial region
are essentially delta functions (perfectly localized), so their contribution to
$\overline{G}_A$ is a multiple of the identity—which is invisible to
$\arctanh$.  The OW surrogate therefore reproduces the exact modular spectrum
to numerical precision regardless of $n_\mathrm{sh}$ or the DW-truncation
flag.

This is a highly non-trivial consistency check: it confirms that the OW
flattened-$\sigma_x$ covariance is in the \emph{same topological phase} and
encodes the \emph{same entanglement pattern} as the exact DW ground state,
at least at the level of the half-system entanglement spectrum.

\subsection{Effect of DW truncation on charge dynamics}

While the spectra are identical, the charge dynamics \emph{does} depend on
which state is used, through the modular eigenvectors $V$.  The DW-truncation
construction forces each OW Wannier mode to have support strictly within the
topological or trivial region.  As a result, $\Hflat$ (and hence
$\overline{G}_A$) becomes exactly block-diagonal in the two-region
decomposition, and the modular eigenstates inherit this block structure.
An excitation injected inside the topological slab therefore has \emph{no
overlap} with trivial-region modular modes and cannot spread across the DW.

This motivates a clear physical picture: DW truncation is equivalent to
imposing a \emph{superselection rule} on the modular dynamics that exactly
decouples the two regions.  The cost is a slight distortion of the modular
eigenvectors relative to the exact ground state (as reflected in the
$\approx 4$--$10\%$ difference in charge profiles), but the topological
structure is preserved.

\subsection{Implications for the adaptive circuit}

The adaptive Markov circuit uses OW projectors as measurement targets.  The
present analysis shows that:
\begin{enumerate}
  \item All OW variants share the exact modular spectrum of the DW ground
    state, confirming they are faithful representations of the topological
    entanglement at the spectral level.
  \item The DW-truncation construction imposes strict topological confinement
    on modular-time dynamics, which is beneficial when one wants to isolate
    topological-region dynamics from bulk fluctuations in the adaptive circuit.
  \item The standard OW variants allow small cross-DW correlations that are
    present in the exact ground state; these may matter when computing
    observables sensitive to the spatial profile of the modular eigenstates
    rather than just their energies.
\end{enumerate}

%% ─────────────────────────────────────────────────────────────────────────────
\section{Computational Details}
\label{sec:numerics}

All simulations use the \texttt{classA\_U1FGTN} class from
\texttt{src/fgtn/classA\_U1FGTN.py}.  The covariance matrix is obtained via
exact diagonalization ($N = 2 \Nx \Ny = 1600$ modes) in double-precision
complex arithmetic.  The modular Hamiltonian is computed via
\texttt{numpy.linalg.eigh} applied to $\overline{G}_A$ (800 modes), followed
by the elementwise $\arctanh$ in \eqref{eq:hmod} with a regularization
$\epsilon = 10^{-10}$ clipping $|g_i|$ away from unity.  The spectral
evolution in \eqref{eq:phit} is performed by projecting $\phi_0$ onto the
modular eigenbasis and applying phase factors $e^{-i\varepsilon_i t}$
analytically, avoiding any time-stepping error.

Output trajectories are stored as compressed NumPy archives
(\texttt{.npz}) in the respective subdirectories.  The full set of 18
trajectories (9 state variants $\times$ 2 injection configs) is recorded in
\texttt{charge\_spreading\_summary.parquet}.  Figure generation uses
\texttt{matplotlib} with the CMU Sans Serif font family at 300\,dpi.

\begin{acknowledgments}
Computation performed on a 112-core shared workstation using
\texttt{numpy}/\texttt{scipy} with CPU thread allocation capped at 14 threads
via \texttt{OMP\_NUM\_THREADS} and related environment variables.
\end{acknowledgments}

\begin{thebibliography}{10}
\bibitem{Peschel2003}
  I.~Peschel,
  \emph{Calculation of reduced density matrices from correlation functions},
  J.~Phys.\ A: Math.\ Gen.\ \textbf{36}, L205 (2003).

\bibitem{Haldane1988}
  F.~D.~M.~Haldane,
  \emph{Model for a quantum Hall effect without Landau levels},
  Phys.~Rev.~Lett.\ \textbf{61}, 2015 (1988).

\bibitem{LiHaldane2008}
  H.~Li and F.~D.~M.~Haldane,
  \emph{Entanglement spectrum as a generalization of entanglement entropy},
  Phys.~Rev.~Lett.\ \textbf{101}, 010504 (2008).

\bibitem{Bisognano1975}
  J.~J.~Bisognano and E.~H.~Wichmann,
  \emph{On the duality condition for a Hermitian scalar field},
  J.~Math.~Phys.\ \textbf{16}, 985 (1975).

\bibitem{Qi2012}
  X.-L.~Qi, H.~Katsura, and A.~W.~W.~Ludwig,
  \emph{General relationship between the entanglement spectrum and the
  edge state spectrum of topological quantum states},
  Phys.~Rev.~Lett.\ \textbf{108}, 196402 (2012).
\end{thebibliography}

\end{document}
